%% Extrait du diagramme des exigences (req) du pont transbordeur.
%% Chaque exigence = une boîte à deux compartiments (stéréotype + nom / Id + Text).
%% Liens : décomposition (cercle barré d'une croix, côté exigence générale) et «refine».
%% Mise en page « téléphone » (largeur < 9 cm) : les exigences sont empilées en
%% deux colonnes de trois rangées, la légende occupant la case libre en bas à gauche.
\documentclass[tikz,border=4pt]{standalone}
\usepackage{liaisons}
\begin{document}
\begin{tikzpicture}[x=1cm,y=1cm,
  cadre/.style={draw=black!70, line width=0.9pt},
  lien/.style={line width=0.8pt, black!75},
  raffine/.style={-{Triangle[length=5pt,width=4.5pt]}, line width=0.8pt, black!75,
                  dash pattern=on 3pt off 2pt},
  boite/.style={draw=solideB, line width=0.9pt, fill=white, text width=3.6cm, inner sep=2.5pt,
                font=\scriptsize,
                execute at begin node={\linespread{0.9}\selectfont
                                       \hyphenpenalty=10000\exhyphenpenalty=10000}},
  tete/.style={boite, align=center},
  corps/.style={boite, align=left}]

%% ---- macro « exigence » : \exig{nom TikZ}{x}{y}{stéréotype}{nom}{Id}{Text} ----
\newcommand\exig[7]{%
  \node[tete, anchor=north] (#1h) at (#2,#3) {\textit{#4}\\[0.5pt] \textbf{#5}};
  \node[corps, anchor=north] (#1b) at (#1h.south) {Id = "#6"\\ Text = "#7"};}
%% ---- macro « cercle barré » (décomposition) : \decompn{nom}{position} ----
\newcommand\decompn[2]{%
  \node[circle, draw=black!75, line width=0.7pt, fill=white, inner sep=0pt,
        minimum size=0.26cm] (#1) at #2 {};
  \draw[lien] (#1.west) -- (#1.east) (#1.south) -- (#1.north);}

%% ---- repères de la grille -------------------------------------------
\def\cgauche{2.10} \def\cdroite{6.35}

%% ================== cadre à onglet ====================================
\draw[cadre] (0,0.35) rectangle (8.45,6.90);
\draw[cadre] (0,6.47) -- (2.35,6.47) -- (2.60,6.70) -- (2.60,6.90);
\node[font=\scriptsize, anchor=west, inner sep=3pt] at (0.1,6.69) {\texttt{req} [Paquet] Exigences};

%% ================== rangée 1 : l'exigence générale et sa décomposition ====
\exig{r1}{\cgauche}{6.30}{«requirement»}{Franchissement de la Loire}{1}%
        {Traverser la Loire à pied ou en busway}
\exig{r11}{\cdroite}{6.30}{«requirement»}{Transport de véhicules et piétons}{1.1}%
        {Charge utile de la nacelle : 100 tonnes}
%% « 1 » se décompose en « 1.1 » : le cercle barré est du côté de l'exigence générale
\decompn{d1}{($(r1h.east)+(0.13,0)$)}
\draw[lien] (d1.east) -- (r11h.west);

%% ================== rangée 2 : les deux exigences fonctionnelles ==========
\exig{r111}{\cgauche}{4.15}{«functional Requirement»}{Traversée à pied}{1.1.1}%
        {Deux cabines : 270 personnes}
\exig{r112}{\cdroite}{4.15}{«functional Requirement»}{Traversée en busway}{1.1.2}%
        {Un busway par traversée}
%% « 1.1 » se décompose en « 1.1.1 » et « 1.1.2 »
\decompn{d11}{($(r11b.south)+(0,-0.13)$)}
\path (d11.south) ++(0,-0.14) coordinate (bus);
\draw[lien] (d11.south) -- (bus);
\draw[lien] (bus) -| (r111h.north);
\draw[lien] (bus) -- (r112h.north);

%% ================== rangée 3 : l'exigence physique (et la légende) ========
\exig{r5}{\cdroite}{2.30}{«physical Requirement»}{Dimensions d'un busway}{5}%
        {Largeur 2,55 m ; hauteur 3,12 m ; 150 personnes}
%% lien «refine» : « 5 » précise l'exigence « 1.1.2 »
\draw[raffine] (r5h.north) -- (r112b.south);
\node[font=\scriptsize, anchor=west, inner sep=2pt] at (6.45,2.52) {«refine»};

%% ================== légende (case libre, en bas à gauche) =================
\decompn{dl}{(0.42,1.98)}
\node[font=\scriptsize, anchor=north west, align=left, inner sep=2pt, text width=3.35cm,
      execute at begin node={\linespread{0.9}\selectfont}] at (0.62,2.11)
     {: l'exigence générale se décompose en exigences plus détaillées};
\node[font=\scriptsize, anchor=north west, align=left, inner sep=2pt, text width=3.9cm,
      execute at begin node={\linespread{0.9}\selectfont}] at (0.20,1.15)
     {«refine» : précise l'exigence pointée};
\end{tikzpicture}
\end{document}
